\usepackage[spanish]{babel}
\addto\captionsspanish{
\def\bibname{Referencias}
\def\tablename{Tabla}
}
\usepackage[hyperfootnotes=false]{hyperref}
\hypersetup{
    colorlinks=true,
    linkcolor=blue,
    filecolor=blue,
    urlcolor=blue,
    citecolor=blue
}
\usepackage{rotating}
\usepackage{pdflscape}
\usepackage{longtable}
\usepackage{tabularx}
\usepackage{array}
\usepackage{enumitem}   % para listas compactas
\usepackage{booktabs}   % opcional, por si quieres líneas más limpias (no imprescindible)
\usepackage{hhline}
%%\usepackage{algorithm}
%\usepackage{algpseudocode}
%\usepackage[noend]{algpseudocode}
\usepackage[nohyperlinks]{acronym}
\usepackage{siunitx}
\usepackage{mathtools}
\usepackage[T1]{fontenc}
\usepackage[utf8]{inputenc}
\usepackage{anysize}
\usepackage{txfonts}
\usepackage{wrapfig} %figura entre texto
\usepackage{times}
\usepackage{listings}
\usepackage{array}
\usepackage{textcomp}
\usepackage{fancyvrb}
\usepackage{fancyhdr}
\usepackage{wallpaper}
\usepackage{verbatim}%Agregado por Luis
\usepackage{latexsym}
\usepackage{url}
\usepackage{multicol}
\usepackage{graphicx}
\usepackage{tikz}
\usepackage{multirow}
\usepackage{float}
\usepackage{lmodern}
\usepackage{eso-pic}
\usepackage{subfig}
\usepackage{longtable}
%\usepackage{amssymb}
\usepackage[dvipsnames]{xcolor}
\usepackage{ragged2e}
\usepackage[all]{xy}
\usepackage{xspace,epic,eepic}
\usepackage{algorithmic}
\usepackage[ruled,vlined]{algorithm2e}
%\usepackage[ruled,vlined,linesnumbered,titlenotnumbered, %portuguese]{algorithm2e}
\usepackage{enumitem}
\usetikzlibrary{positioning,arrows.meta,svg.path}
\setlist{topsep=0pt,noitemsep}
\setcounter{tocdepth}{3}
\marginsize{3cm}{2.5cm}{2.5cm}{2.5cm}

%Definicion de Colores
\definecolor{gray97}{gray}{.97}
\definecolor{gray75}{gray}{.75}
\definecolor{gray45}{gray}{.45}
\definecolor{listinggray}{gray}{0.9}
\definecolor{lbcolor}{rgb}{0.9,0.9,0.9}
\newcommand\rojo[1]{\textcolor[rgb]{1,0,0}{\textbf{#1}}}
\newcommand\red[1]{\textcolor[rgb]{1.00,0.00,0.00}{#1}}
\newcommand\azul[1]{\textcolor[rgb]{0,0,1}{\textbf{#1}}}
\newcommand\blue[1]{\textcolor[rgb]{0,0,1}{{#1}}}
\newcommand\verde[1]{\textcolor[rgb]{0,.5,0.2}{\textbf{#1}}}
\newcommand\naranjo[1]{\textcolor[rgb]{1.00,0.36,0.06}{\textbf{#1}}}
%\ifCLASSINFOpdf
%\else
%\fi
%\hyphenation{pa-la-bra}

\lstset{
	%backgroundcolor=\color{lbcolor},
	tabsize=4,
	rulecolor=,
	language=SQL,
        basicstyle=\small,
        upquote=true,
        aboveskip={1.5\baselineskip},
        columns=fixed,
        showstringspaces=false,
        extendedchars=true,
        breaklines=true,
        prebreak = \raisebox{0ex}[0ex][0ex]{\ensuremath{\hookleftarrow}},
        %frame=single,
        showtabs=false,
        showspaces=false,
        showstringspaces=false,
        identifierstyle=\ttfamily,
        keywordstyle=\color[rgb]{0,0,0},
        commentstyle=\color[rgb]{0,0,0},
        stringstyle=\color[rgb]{0,0,0},
}

\lstdefinestyle{javadto}{
  language=Java,
  basicstyle=\small\ttfamily,
  keywordstyle=\color[HTML]{7B30D0}\bfseries,
  stringstyle=\color[HTML]{067D17},
  commentstyle=\color[HTML]{6A737D},
  emphstyle=\color[HTML]{0033B3}\bfseries,
  emph={String,Long,int,boolean,BigDecimal,List,LocalDateTime,MultipartFile},
  showstringspaces=false,
  frame=single,
  framerule=0.4pt,
  rulecolor=\color[HTML]{CCCCCC},
  backgroundcolor=\color[HTML]{F7F7F7},
  aboveskip=8pt,
  belowskip=8pt,
  xleftmargin=4pt,
  xrightmargin=4pt,
  tabsize=4,
  breaklines=true,
}

%\renewcommand{\labelenumi}{\arabic{enumi}.} % (1., 2., 3.,...)
%\renewcommand{\labelenumi}{\roman{enumi}.} %  (i., ii., iii.,...)
%\renewcommand{\labelenumi}{\Roman{enumi}.} %  (I., II., III.,...)
%\renewcommand{\labelenumi}{\alph{enumi}.}   % (a., b., c.,...)
\renewcommand{\labelenumi}{(\alph{enumi})} % [(a), (b), (c),...]
%\renewcommand{\labelenumi}{\Alph{enumi}.}  %  (A., B., C.,...)


\newtheorem{theorem}{Theorem}[chapter]
\newtheorem{teorema}[theorem]{Teorema}
\newtheorem{lem}{Theorem}[chapter]
\newtheorem{lema}[lem]{Lema}
\newtheorem{prop}{Theorem}[chapter]
\newtheorem{proposition}[prop]{Proposición}
\newtheorem{coro}{Theorem}[chapter]
\newtheorem{corolario}[coro]{Corolario}
\newtheorem{exam}{Theorem}[chapter]
\newtheorem{example}[exam]{Ejemplo}
\newtheorem{test}{Theorem}[chapter]
\newtheorem{prueba}[test]{Prueba}
\newtheorem{defi}{Theorem}[chapter]
\newtheorem{definition}[defi]{Definición}
\newtheorem{obs}{Theorem}[chapter]
\newtheorem{observacion}[obs]{Observación}

\newcommand{\keywords}[1]{\par\addvspace\baselineskip
\noindent\keywordname\enspace\ignorespaces#1}
\renewcommand{\abstractname}{\prefacesection{Resumen}}
%\renewcommand{\tableofcontents}{Índice General}
\renewcommand{\listtablename}{Índice de Tablas}
\renewcommand{\tablename}{Tabla}
\renewcommand{\refname}{Bibliografía}
\newcommand{\boxtheorem}{\hfill $\Box$}
%%%%%%%%IEEE Palabras Claves
%\renewcommand{\IEEEkeywords}{\textbf{\emph{Palabras Clave---}}}
